\scanpage{127}

\uline{Lemme} \deleted{2}3 \quad Soit $f:X\unsure{\text{---}|\longrightarrow}Y$ un morphisme \inserted{compactifiable} de type fini, de dim relative $\leq d'$, $F$ un $\mathbb Z_\ell$-faisceau admissible de type $\leq d$ sur $X$, alors pour \deleted{\unsure{tt|tout}} $i$, $R^i_!f(F)$ sur $Y$ sont admissibles de type $\leq d+d'$

\margin{$\Big\downarrow$}

On peut supposer $Y=\operatorname{Sp}k$, $k$ un corps.

Récurrence sur $d+d'$, trivial si $d'\leq0$. Supposons $d+d'>0$, et le th. démontré pour les \deleted{\unsure{$d'$}} \inserted{$\delta'<d'$}.

Soit $U\subset X$ \inserted{ouvert} lisse sur $K$ et dense dans $X$, et \inserted{que $F|U$ ait une filtration $d$-admissible.} [NB. on est obligé dans la déf. plus haut d'admettre une extension radicielle sur la base.] et soit $Z=X-U$.

On a

\[ \unsure{\mathbb R|R}^{\unsure{-|i}}_! f_U(F|U)\longrightarrow \mathbb R^{-}_!f_*(F)\unsure{\text{---}|\longrightarrow}R^{-}_!f_{\unsure{Z|Y}*}(F|Y) \]

ce qui ns ramène au cas de $f_Z$ et $f_U$ grâce aux lemmes 1 et 2 ; le cas de $f_Z$ est contenu dans l'hyp. de réc, et \deleted{\unsure{alors}} \unsure{pour|donc} le cas \deleted{\unsure{de}} où $X$ est lisse sur $k$ et a une filtration $d$-admissible.

Par dévissage et le lemme 1 \deleted{\unsure{de}} on peut supposer que $F$ est un quotient \inserted{ess. const.} d'un

% archive sheet number: 126
